\section{Related work}

\subsection{State estimation and complementary filtering}

The Kalman filter established recursive prediction and measurement correction
for linear stochastic systems, with gains determined by propagated covariance
\cite{kalman1960_filtering}. Complementary filters provide a lower-complexity
alternative when measurements and model predictions have different reliability
or frequency characteristics. Nonlinear complementary observers on
$\mathrm{SO}(3)$ provide an influential example for attitude estimation
\cite{mahony2008_complementary}. Extended-state observers likewise estimate
unmodeled interactions from motion measurements; recent IJDC work demonstrates
this role for wheel--rail dynamics \cite{zhao2023_wheel_rail_observer}.

Learning has been introduced into this estimation loop in several ways.
KalmanNet learns Kalman-style updates for partially known dynamics
\cite{revach2022_kalmannet}; autodifferentiable ensemble Kalman filters expose
the assimilation procedure to gradient-based learning
\cite{chen2022_autodifferentiable_enkf}; hybrid estimators combine physical and data-driven
components in energy and vehicle systems
\cite{chen2022_hybrid_physics_data_driven_online,chen2023_hybrid_physics_and_data_driven};
and machine learning with data assimilation connects estimation to uncertainty
quantification \cite{cheng2023_machine_learning_with_data_assimilation}. CFR-WM
uses a fixed-gain mechanism: an observed channel receives a prespecified scalar
gain, a missing channel receives zero gain, and the corrected state initializes
a learned multistep predictor. A classical fixed-gain-filter-plus-nominal-
rollout baseline separates the context filter from the neural residual.

\subsection{Physics-guided dynamical learning and predictive control}

Physics-guided learning can constrain outputs, embed equations, combine
numerical solvers with neural corrections, or assimilate uncertain
observations. Reviews emphasize that the appropriate integration depends on the
dynamics, available measurements, and downstream objective
\cite{cuomo2022_scientific_machine_learning_through_physics,bradley2022_perspectives_on_the_integration_between,yu2024_learning_dynamical_systems_from_data}.
Differentiable modeling exposes mechanistic components to gradient-based
learning, while horizon-aware optimization addresses temporal error
accumulation
\cite{shen2023_differentiable_modelling_to_unify_machine,wang2024_respecting_causality_for_training_physics,list2022_learned_turbulence_modelling_with_differentiable}.

Physics-guided neural networks have been evaluated on nonlinear dynamics
\cite{robinson2022_physics_guided_neural_networks_for}, and physics-guided Deep
Markov Models combine state-space physics with probabilistic latent inference
\cite{liu2022_physics_guided_deep_markov_models}. These approaches primarily
shape the transition or latent dynamics. CFR-WM instead couples a fixed context
correction to a bounded residual transition, allowing the forecast-origin and
rollout contributions to be evaluated separately. This distinction also
matters for learning-based model predictive control, where a learned prediction
model is repeatedly queried across an optimization horizon and uncertainty or
model mismatch can affect feasibility and performance
\cite{hewing2020_learning_mpc}. Sparse system identification has likewise been
coupled directly to nonlinear model predictive control in the low-data regime
\cite{kaiser2018_sindy_mpc}. The present study focuses on the offline
estimator--predictor mapping that supplies state forecasts to such controllers.

\subsection{World models and predictive uncertainty}

World models use recurrent stochastic state-space models, normalized latent
representations, transformers, and other sequence architectures. DreamerV3 and
TD-MPC2 demonstrate scalable learned dynamics for continuous control, while
IRIS and STORM show transformer-based alternatives
\cite{hafner2025_mastering_diverse_control_tasks_through,hansen2023_td_mpc2_scalable_robust_world,micheli2022_transformers_are_sample_efficient_world,zhang2023_storm_efficient_stochastic_transformer_based}.
For the explicit-state forecasting problem here, we implement compact,
approximately parameter-matched recurrent, normalized, transformer, and
physics-guided transition families under a common protocol. Broader surveys
organize world models around representation, dynamics, planning, and agent
learning \cite{ding2025_understanding_world_or_predicting_future}.

Predictive uncertainty can represent observation noise, transition ambiguity,
or model-form error, and likelihood alone does not establish calibrated
intervals \cite{abdar2021_a_review_of_uncertainty_quantification}. Dynamic-
system studies assess coverage and interval width together, while Bayesian and
conformal methods address model-form and distributional uncertainty
\cite{klotz2022_uncertainty_estimation_with_deep_learning,ramancha2022_accounting_for_model_form_uncertainty,portela2025_conformal_prediction_for_uncertainty_quantification}.
We therefore report Gaussian negative log likelihood, empirical coverage,
interval width, and calibration error in addition to multistep point accuracy.
